% c04 — source pages 32–41, questions 170–224 (grammatical sentences, linkers, compound words, prefixes and suffixes)
\begin{questions}

\Q{170} వ్యాకరణపరంగా సరైన వాక్యాన్ని గుర్తించండి.
\opts{\en{The children have listening to stories.}}{\en{The children was listening to stories.}}{\en{The children are listening to stories.}}{\en{The children has listening to stories.}}

\Q{171} వ్యాకరణపరంగా సరైన వాక్యాన్ని గుర్తించండి.
\opts{\en{Stephen is one of the best singers.}}{\en{Stephen are one of the best singers.}}{\en{Stephen have one of the best singers.}}{\en{Stephen will one of the best singers.}}

\Q{172} వ్యాకరణపరంగా సరైన వాక్యాన్ని గుర్తించండి.
\opts{\en{They were sold their bungalow.}}{\en{They are sold their bungalow.}}{\en{They has sold their bungalow.}}{\en{They have sold their bungalow.}}

\Q{173} వ్యాకరణపరంగా సరైన వాక్యాన్ని గుర్తించండి.
\opts{\en{The crow is fly high.}}{\en{The crow has fly high.}}{\en{The crow is flying high.}}{\en{The crow will flying high.}}

\Q{174} వ్యాకరణపరంగా సరైన వాక్యాన్ని గుర్తించండి.
\opts{\en{My uncle has bought a car.}}{\en{My uncle has buying a car.}}{\en{My uncle have bought a car.}}{\en{My uncle have buying a car.}}

\Q{175} వ్యాకరణపరంగా సరైన వాక్యాన్ని గుర్తించండి.
\opts{\en{John is going to school.}}{\en{John are going to school.}}{\en{John have going to school.}}{\en{John had going to school.}}

\Q{176} వ్యాకరణపరంగా సరైన వాక్యాన్ని గుర్తించండి.
\opts{\en{Rajesh was a billionaire.}}{\en{Rajesh have a billionaire.}}{\en{Rajesh have been a billionaire.}}{\en{Rajesh will a billionaire.}}

\Q{177} వ్యాకరణపరంగా సరైన వాక్యాన్ని గుర్తించండి.
\opts{\en{You has very naughty in childhood.}}{\en{You have very naughty in childhood.}}{\en{You had very naughty in childhood.}}{\en{You were very naughty in childhood.}}

\Q{178} వ్యాకరణపరంగా సరైన వాక్యాన్ని గుర్తించండి.
\opts{\en{All the servants goes to their homes.}}{\en{All the servants going to their homes.}}{\en{All the servants went to their homes.}}{\en{All the servants will going to their homes.}}

\Q{179} వ్యాకరణపరంగా సరైన వాక్యాన్ని గుర్తించండి.
\opts{\en{The thief cutting his stick by two inches.}}{\en{The thief has cutting his stick by two inches.}}{\en{The thief have cutting his stick by two inches.}}{\en{The thief cut his stick by two inches.}}

\Q{180} వ్యాకరణపరంగా సరైన వాక్యాన్ని గుర్తించండి.
\opts{\en{Once there was a merchant.}}{\en{Once there has a merchant.}}{\en{Once there have a merchant.}}{\en{Once there were a merchant.}}

\Q{181} వ్యాకరణపరంగా సరైన వాక్యాన్ని గుర్తించండి.
\opts{\en{The thief was one of his servants.}}{\en{The thief were one of his servants.}}{\en{The thief have been one of his servants}}{\en{The thief will one of his servants.}}

\Q{182} వ్యాకరణపరంగా సరైన వాక్యాన్ని గుర్తించండి.
\opts{\en{His teeth has shining.}}{\en{His teeth are shining.}}{\en{His teeth have shining.}}{\en{His teeth was shining.}}

\Q{183} వ్యాకరణపరంగా సరైన వాక్యాన్ని గుర్తించండి.
\opts{\en{Neither the chairman nor the directors are present.}}{\en{Neither the chairman nor the directors was present.}}{\en{Neither the chairman nor the directors has present.}}{\en{Neither the chairman nor the directors have present.}}

\Q{184} వ్యాకరణపరంగా సరైన వాక్యాన్ని గుర్తించండి.
\opts{\en{One of the servants have his stick shorter.}}{\en{One of the servants have been his stick shorter.}}{\en{One of the servants had his stick shorter.}}{\en{One of the servants were his stick shorter.}}

\Q{185} వ్యాకరణపరంగా సరైన వాక్యాన్ని గుర్తించండి.
\opts{\en{The quality of mangoes were not good.}}{\en{The quality of mangoes was not good.}}{\en{The quality of mangoes have not good.}}{\en{The quality of mangoes has not good.}}

\Q{186} వ్యాకరణపరంగా సరైన వాక్యాన్ని గుర్తించండి.
\opts{\en{The novelist is dead.}}{\en{The novelist are dead.}}{\en{The novelist were dead.}}{\en{The novelist will dead.}}

\Q{187} వ్యాకరణపరంగా సరైన వాక్యాన్ని గుర్తించండి.
\opts{\en{The house with its contents was insured.}}{\en{The house with its contents were insured.}}{\en{The house with its contents has insured.}}{\en{The house with its contents have insured.}}

\Q{188} వ్యాకరణపరంగా సరైన వాక్యాన్ని గుర్తించండి.
\opts{\en{Silver as well as cotton has fallen in price.}}{\en{Silver as well as cotton have fallen in price.}}{\en{Silver as well as cotton have been fallen in price.}}{\en{Silver as well as cotton fallen in price.}}

\Q{189} వ్యాకరణపరంగా సరైన వాక్యాన్ని గుర్తించండి.
\opts{\en{Neither food nor water was to be found there.}}{\en{Neither food nor water were to be found there.}}{\en{Neither food nor water have to be found there.}}{\en{Neither food nor water are to be found there.}}

\Q{190} సరైన అనుసంధాన పదాన్ని \en{(Linker)} ఎంచుకోండి.\par
నేను గిన్నెను కింద పడేశాను, \blank\ అది పగలలేదు.
\eng{I dropped the dish, \blank\ it did not break.}
\opts{\en{so}}{\en{as}}{\en{but}}{\en{though}}

\Q{191} కింది వాక్యాన్ని పూర్తి చేయడానికి సరైన అనుసంధాన పదాన్ని \en{(Linker)} ఎంచుకోండి.\par
నువ్వు \blank\ నేను వెళ్ళను.
\eng{I shall not go \blank\ you do.}
\opts{\en{as soon as}}{\en{scarcely}}{\en{unless}}{\en{then}}

\Q{192} సరైన అనుసంధాన పదాన్ని \en{(Linker)} ఎంచుకోండి.\par
\blank\ నాకు ఆ దుస్తులు నచ్చాయి, నేను వాటిని కొనకూడదని నిర్ణయించుకున్నాను.
\eng{\blank\ I liked the dress, I decided not to buy it.}
\opts{\en{But}}{\en{Though}}{\en{Because}}{\en{As soon as}}

\Q{193} సరైన అనుసంధాన పదాన్ని \en{(Linker)} ఎంచుకోండి.\par
మా గది చాలా చిన్నది, \blank\ మేము నిజంగా పట్టించుకోలేదు.
\eng{Our room was very small, \blank\ we did not really mind}
\opts{\en{if}}{\en{because}}{\en{due to}}{\en{but}}

\Q{194} సరైన అనుసంధాన పదాన్ని \en{(Linker)} ఎంచుకోండి.\par
మేము బయటకు భోజనానికి వెళ్ళకూడదని నిర్ణయించుకున్నాం \blank\ మేము మరీ అలసిపోయాం.
\eng{We decided not to go out for a meal \blank\ we were simply too tired.}
\opts{\en{because}}{\en{so}}{\en{but}}{\en{though}}

\Q{195} కింది వాక్యాన్ని పూర్తి చేయడానికి సరైన ఐచ్ఛికాన్ని ఎంచుకోండి.\par
నువ్వు మర్చిపోవడానికి \blank\ దాన్ని చేసేయ్.
\eng{Do it \blank\ you forget.}
\opts{\en{beforehand}}{\en{after}}{\en{before}}{\en{yet}}

\Q{196} కింది వాక్యాన్ని పూర్తి చేయడానికి సరైన ఐచ్ఛికాన్ని ఎంచుకోండి.\par
వారి ఇల్లు \blank\ పెద్దది \blank\ చిన్నది.
\eng{Their house is \blank\ big \blank\ small.}
\opts{\en{neither, nor}}{\en{either, nor}}{\en{neither, not}}{\en{neither, or}}

\Q{197} సరైన అనుసంధాన పదాన్ని \en{(Linker)} ఎంచుకోండి.\par
ఒక ప్రమాదం \blank\ రోడ్డు మూసివేయబడింది.
\eng{The road was closed \blank\ an accident.}
\opts{\en{due to}}{\en{in spite of}}{\en{unless}}{\en{furthermore}}

\Q{198} సరైన అనుసంధాన పదాన్ని \en{(Linker)} ఎంచుకోండి.\par
ఆ దుస్తులు సరిపోలేదు, \blank\ ప్రకృతి వాటిని తిరిగి ఇచ్చేయవలసి వచ్చింది.
\eng{The dress didn’t fit, \blank\ Prakruthi had to return it.}
\opts{\en{but}}{\en{if}}{\en{or}}{\en{so}}

\Q{199} తగిన అనుసంధాన పదాన్ని \en{(Linker)} ఎంచుకోండి.\par
కేఫ్ రద్దీగా ఉంది, \blank\ మాకు ఒక టేబుల్ దొరికింది.
\eng{The café was crowded, \blank\ we found a table.}
\opts{\en{and}}{\en{so}}{\en{if}}{\en{but}}

\Q{200} కింది వాక్యాన్ని పూర్తి చేయడానికి సరైన ఐచ్ఛికాన్ని ఎంచుకోండి.\par
ఆమె తెలివైనది మాత్రమే కాదు, \blank\ సంగీతంలో ఎంతో ప్రతిభ గలది.
\eng{She was not only intelligent \blank\ very musical.}
\opts{\en{also}}{\en{but also}}{\en{if}}{\en{neither}}

\Q{201} కింది వాక్యాన్ని పూర్తి చేయడానికి సరైన ఐచ్ఛికాన్ని ఎంచుకోండి.\par
నడవాకు \blank\ చివర ఒక తలుపు ఉంది.
\eng{There is a door at \blank\ end of the corridor.}
\opts{\en{nor}}{\en{or}}{\en{not}}{\en{either}}

\Q{202} సరైన అనుసంధాన పదాన్ని \en{(Linker)} ఎంచుకోండి.\par
మా అక్కకు ఉద్యోగం వచ్చింది, \blank\ ఆమె దానిని ఊహించలేదు.
\eng{My sister got the job, \blank\ she didn’t expect to.}
\opts{\en{but}}{\en{or}}{\en{in spite of}}{\en{so}}

\Q{203} సరైన అనుసంధాన పదాన్ని \en{(Linker)} ఎంచుకోండి.\par
నేరాలను తగ్గించే ఉద్దేశంతో ప్రభుత్వం ఈ చర్యలు తీసుకుంది.
\eng{The government took these measures \blank\ reduce crime.}
\opts{\en{in spite of}}{\en{in case of}}{\en{in order to}}{\en{neither nor}}

\Q{204} తగిన అనుసంధాన పదాన్ని \en{(Linker)} ఎంచుకోండి.\par
కొత్త వ్యవస్థ మరింత సమర్థవంతంగా ఉంటుందని భావించారు. \blank, ఆచరణలో అది గందరగోళాన్ని సృష్టించింది.
\eng{The new system was supposed to be more efficient. \blank, in practice it caused chaos.}
\opts{\en{However}}{\en{And}}{\en{If}}{\en{In spite of}}

\Q{205} కింది వాక్యాలను పూర్తి చేయడానికి సరైన అనుసంధాన పదాలను \en{(Linkers)} ఎంచుకోండి.\par
\blank, అలారం నన్ను నిద్రలేపుతుంది. \blank\ నేను స్నానం చేస్తాను.
\eng{\blank, the alarm wakes me up. \blank\ I take bath.}
\opts{\en{First, Then}}{\en{Then, First}}{\en{After that, First}}{\en{Finally, First}}

\Q{206} సరైన అనుసంధాన పదాన్ని \en{(Linker)} ఎంచుకోండి.\par
రోడ్డు నీట మునిగింది. \blank\ పోలీసులు దానిపై రాకపోకలను నిలిపివేశారు.
\eng{The road was under water. The police \blank\ closed it to traffic.}
\opts{\en{as}}{\en{therefore}}{\en{but}}{\en{however}}

\Q{207} సరైన అనుసంధాన పదాన్ని \en{(Linker)} ఎంచుకోండి.\par
అది ఒక భయంకరమైన ప్రయాణం. \blank, చివరికి మేము సురక్షితంగా అక్కడికి చేరుకున్నాం.
\eng{It was a terrible journey. \blank, we got there safely in the end.}
\opts{\en{Besides}}{\en{Consequently}}{\en{Therefore}}{\en{Still}}

\Q{208} సరైన అనుసంధాన పదాన్ని \en{(Linker)} ఎంచుకోండి.\par
విమాన రాకపోకల నియంత్రణాధికారుల సమ్మె మొదలైంది. \blank\ చాలా విమానాలు రద్దయ్యాయి.
\eng{A strike by air traffic controllers has begun. Many flights have \blank\ been cancelled.}
\opts{\en{besides}}{\en{consequently}}{\en{in spite of}}{\en{but}}

\Q{209} సరైన అనుసంధాన పదాన్ని \en{(Linker)} ఎంచుకోండి.\par
ఆమె బూట్లు వేసుకుంది \blank\ తన పాదాలు తడవవు.
\eng{She wore boots \blank\ her feet would not get wet.}
\opts{\en{so that}}{\en{if}}{\en{but}}{\en{to}}

\Q{210} కింది వాక్యాన్ని పూర్తి చేయడానికి సరైన ఐచ్ఛికాన్ని ఎంచుకోండి.\par
తాను ఆమెతో ఆ విషయం గురించి చర్చించలేదని అతను చెప్పాడు. \blank, అతను ఆమెను కనీసం సంప్రదించను కూడా లేదు.
\eng{He said he had not discussed the matter with her. \blank, he had not even contacted her.}
\opts{\en{farther}}{\en{for other more}}{\en{furthermore}}{\en{farther more}}

\Q{211} కింది వాటి నుండి సమాస పదాలను \en{(Compound words)} ఎంచుకోండి.\par
\en{a.} దిష్టిబొమ్మ \en{(scarecrow)}\par
\en{b.} నల్లబల్ల \en{(blackboard)}\par
\en{c.} వాయిదా వేయడం \en{(postpone)}\par
\en{d.} పుస్తకాల అర \en{(bookshelf)}\par
సరైన ఐచ్ఛికాన్ని ఎంచుకోండి.
\opts{\en{a, b, c}}{\en{b, c, d}}{\en{a, b, d}}{\en{a, c, d}}

\Q{212} కింది వాటి నుండి సమాస పదాలను \en{(Compound words)} ఎంచుకోండి.\par
\en{a.} ఎవరైనా \en{(anyone)}\par
\en{b.} ఆగు \en{(hang on)}\par
\en{c.} చేతి సంచి \en{(handbag)}\par
\en{d.} ఇంద్రధనుస్సు \en{(rainbow)}\par
సరైన ఐచ్ఛికాన్ని ఎంచుకోండి.
\opts{\en{a, b, c}}{\en{a, c, d}}{\en{b, c, d}}{\en{a, b, c, d}}

\Q{213} కింది వాటి నుండి సమాస పదాలను \en{(Compound words)} ఎంచుకోండి.\par
\en{a.} పొద్దుతిరుగుడు పువ్వు \en{(sunflower)}\par
\en{b.} చక్రాల కుర్చీ \en{(wheelchair)}\par
\en{c.} వార్తాపత్రిక \en{(newspaper)}\par
\en{d.} పారిజాత పువ్వు \en{(shiuli flower)}\par
సరైన ఐచ్ఛికాన్ని ఎంచుకోండి.
\opts{\en{a, b, c}}{\en{c, b, d}}{\en{a, c ,d}}{\en{a, b, c, d}}

\Q{214} కింది వాటి నుండి సమాస పదాలను \en{(Compound words)} ఎంచుకోండి.\par
\en{a.} పర్వత పాదాల వద్ద ఉండే కొండలు \en{(foothills)}\par
\en{b.} కాళ్లకు వేసే ఆధార పట్టీలు \en{(callipers)}\par
\en{c.} పొద్దుతిరుగుడు పువ్వు \en{(sunflower)}\par
\en{d.} ఇంద్రధనుస్సు \en{(rainbow)}\par
సరైన ఐచ్ఛికాన్ని ఎంచుకోండి.
\opts{\en{a, b}}{\en{a, c, d}}{\en{b, c, d}}{\en{a, b, c, d}}

\Q{215} అర్థవంతమైన పదాన్ని ఏర్పరచడానికి సరైన ప్రత్యయాన్ని \en{(Suffix)} ఉపయోగించిన పదాన్ని కనుగొనండి.
\opts{\en{creativeness}}{\en{bitterment}}{\en{honestful}}{\en{cheerfulness}}

\Q{216} వ్యతిరేకార్థం రావడానికి \en{‘un’} ను ఉపసర్గగా \en{(Prefix)} తీసుకోని పదం కింది వాటిలో ఏది?
\opts{నిజమైన \en{(true)}}{శిక్షణ పొందిన \en{(trained)}}{ముఖ్యమైన \en{(important)}}{అర్థం చేసుకోవడం \en{(understand)}}

\Q{217} ‘ఉపసర్గ + మూల పదం’ \en{(Prefix + base word)} సరైన కలయికను ఉపయోగించిన వాక్యాన్ని ఎంచుకోండి.
\opts{\en{She is very \underline{impolite}.}}{\en{Nothing is \underline{unpossible}.}}{\en{The project is \underline{uncompleted}.}}{\en{They always \underline{inobey} the rules.}}

\Q{218} అర్థవంతమైన పదం రావడానికి \en{‘ness’} అనే ప్రత్యయాన్ని \en{(Suffix)} కింది మూల పదాలలో \en{(Base word)} దేనికి చేర్చవచ్చు?
\opts{ఎత్తైన \en{(high)}}{సాధారణ \en{(general)}}{కదులు \en{(move)}}{ఉదారమైన \en{(generous)}}

\Q{219} ఇచ్చిన వాక్యంలోని మూల పదానికి \en{(Base word)} చేర్చి, దానిని వ్యతిరేకార్థకంగా మార్చగల ఉపసర్గను \en{(Prefix)} గుర్తించండి.\par
‘దయచేసి తలుపుకు \tw{తాళం వేయండి}{lock}.’
\eng{‘Please, \underline{lock} the door.’}
\opts{\en{dis}}{\en{un}}{\en{in}}{\en{im}}

\Q{220} ఒకదానికొకటి వ్యతిరేకార్థం గల సరైన పద జంటను ఎంచుకోండి.
\opts{\en{appear-disappear}}{\en{direct-undirect}}{\en{healthy-inhealthy}}{\en{mortal-unmortal}}

\Q{221} కింది పద జంటలలో దేనిలో ఉపసర్గ \en{(Prefix)} మూల పదానికి \en{(Base word)} సరిపోలలేదు?
\opts{\en{legal-illegal}}{\en{correct-incorrect}}{\en{mortal-immortal}}{\en{mature-unmature}}

\Q{222} తగిన పదాలతో ఖాళీలను పూరించండి.\par
రాజు ఆహారాన్ని \blank, అది \blank\ గా ఉంది.
\eng{Raju \blank\ the food and it was \blank.}
\opts{\en{taste, tasted}}{\en{tasted, taste}}{\en{tasted, tasty}}{\en{tasty, tasted}}

\Q{223} తగిన పదంతో ఖాళీని పూరించండి.\par
ఆ లోయ \blank\ కొండలతో, చెట్లతో చుట్టుముట్టబడి ఉంది.
\eng{The valley was surrounded by \blank\ hills and trees.}
\opts{\en{beauty}}{\en{beautiful}}{\en{beautifully}}{\en{beautiness}}

\Q{224} బ్రాకెట్‌లో ఇచ్చిన పదానికి తగిన ప్రత్యయాన్ని \en{(Suffix)} చేర్చండి.\par
\blank\ \en{(Final)}, రాజే స్వయంగా త్రాసులో కూర్చున్నాడు.
\eng{\blank\ (Final), the king himself sat in the scale.}
\opts{\en{y}}{\en{ly}}{\en{lly}}{\en{ily}}

\end{questions}
